\section{Case studies}
\label{app:cases}
% DRAFT (Claude): awaits the author's rewrite. Figure and table are written by
% `markov-lean cases --run results/pilot3 --out paper/figures`.

Figure~\ref{fig:cases} and Table~\ref{tab:cases} show a few short proofs from
the judged sample that clearly move between areas. A step's area is the
\textsc{msc} code both judges gave its goal; it counts as confident when the
judges agree and the lower of their two confidences is at least $0.7$. A
confident switch is a pair of adjacent confident steps with different areas,
and the tactic of the first step is the one that moved the goal. We keep proofs
of 4 to 14 steps whose steps were all labelled, whose tactics are short enough
to read, and on which the judges agree for at least 80\% of the steps; we rank
them by confident switches, then agreement, then length, and where the sample
allows take one proof per Mathlib file and only proofs that add a new pair of
areas. The proofs are chosen to be readable, not typical; the
switch rates in Section~\ref{sec:experiments} describe the typical proof.

\begin{figure*}[t]
  \centering
  \maybefigure{case_studies}{\textwidth}
  \caption{Case-study proofs, one row per proof and one cell per step, coloured
  by the judges' area for the goal before the step. Pale cells are steps where
  the judges agree with low confidence; gray cells marked ``?'' are steps where
  they disagree. Outlined cells start a confident switch, and their tactic,
  shown by its first word under the cell, is in bold.}
  \label{fig:cases}
\end{figure*}

\begin{table}[h]
  \centering
  \caption{The case-study proofs of Figure~\ref{fig:cases}: theorem, Mathlib
  file, the sequence of confident areas, and each switching tactic with its
  step.}
  \label{tab:cases}
  \small
  \maybetable{case_studies_table}
\end{table}
